\documentclass[opre]{informs3} % current default for manuscript submission
\OneAndAHalfSpacedXI % current default line spacing

\usepackage[utf8]{inputenc}
\usepackage{graphicx} % Required for inserting images

\usepackage{bbm}
\usepackage{amsthm}
\usepackage{thmtools}
\usepackage{thm-restate}
\usepackage{endnotes}
\usepackage{amsmath, amssymb, amsthm, mathrsfs}
\usepackage{geometry}
\usepackage{bbm}
%\let\footnote=\endnote
\let\enotesize=\normalsize
\def\notesname{Endnotes}%
\def\makeenmark{$^{\theenmark}$}
\def\enoteformat{\rightskip0pt\leftskip0pt\parindent=1.75em
  \leavevmode\llap{\theenmark.\enskip}}

% Private macros here (check that there is no clash with the style)

% Natbib setup for author-year style
\usepackage{natbib}
 \bibpunct[, ]{(}{)}{,}{a}{}{,}%
 \def\bibfont{\small}%
 \def\bibsep{\smallskipamount}%
 \def\bibhang{24pt}%
 \def\newblock{\ }%
 \def\BIBand{and}%

\usepackage[hyperfootnotes=false, colorlinks=true, linkcolor=black, urlcolor=black, citecolor={blue}]{hyperref}
\usepackage{cleveref}

%% Setup of theorem styles. Outcomment only one.
%% Preferred default is the first option.
\TheoremsNumberedThrough     % Preferred (Theorem 1, Lemma 1, Theorem 2)
%\TheoremsNumberedByChapter  % (Theorem 1.1, Lema 1.1, Theorem 1.2)
\ECRepeatTheorems

%% Setup of the equation numbering system. Outcomment only one.
%% Preferred default is the first option.
\EquationsNumberedThrough    % Default: (1), (2), ...
%\EquationsNumberedBySection % (1.1), (1.2), ...
\allowdisplaybreaks
%\mathtoolsset{showonlyrefs}
\usepackage{mathtools}

\usepackage{bm,bbm,algorithmic,algorithm,bigstrut,color,framed,multirow,tikz,graphicx,enumitem,subcaption,booktabs,wrapfig}
\usetikzlibrary{arrows,automata,shapes.geometric,mindmap, shadows}
\usepackage{listings}
\usepackage{xcolor}
\usepackage[skip=2pt,font=scriptsize]{caption}

\tikzstyle{startstop} = [rectangle, rounded corners, 
minimum width=3cm, 
minimum height=1cm,
text width=4cm,
text centered, 
draw=black, 
fill=blue!30]
\tikzstyle{startstop2} = [rectangle, rounded corners, 
minimum width=3cm, 
minimum height=1cm,
text width=4cm,
text centered, 
draw=black, 
fill=yellow!30]
\tikzstyle{arrow} = [thick,->,>=stealth]


\lstset{
  language=Python,
  basicstyle=\ttfamily\small,
  backgroundcolor=\color{gray!10},
  frame=single,
  breaklines=true,          % Rompe líneas solo si es necesario
  breakatwhitespace=false,  % No cortar palabras automáticamente
  showstringspaces=false,
  tabsize=2,
  captionpos=b,
  keepspaces=true,
  columns=fullflexible      % <- Esta línea es la clave para evitar los saltos raros
}

\newcommand{\appendix}{%
  \par
  \setcounter{section}{0}%
  \setcounter{subsection}{0}%
  \renewcommand{\thesection}{\Alph{section}}%
}



\newcommand{\E}{\mathbb{E}}
\newcommand{\Pb}{\mathbb{P}}
\newcommand{\Real}{\mathbb{R}}
\newcommand{\Natu}{\mathbb{N}}
\newcommand{\Int}{\mathbb{Z}}
\newcommand{\set}[1]{\left\{#1\right\}}
\newcommand{\abs}[1]{\left|#1\right|}
\newcommand{\bra}[1]{\left(#1\right)}
\newcommand{\tran}[1]{\left(#1\right)^\top }
\newcommand{\infnorm}[1]{\|#1\|_{\infty}}
\newcommand{\setnorm}[2]{\|#1\|_{#2}}
\newcommand{\conc}[2]{\left[\begin{array}{c}#1\\#2\end{array}\right]}
% \newcommand{\ind}[1]{\mathbb{1}\set{#1}}
\newcommand{\ind}[1]{\mathbbm{1}_{#1}}
\newcommand{\Ind}[1]{\mathbbm{1}\{#1\}}
\newcommand{\crs}[1]{\left\langle#1\right\rangle}
\newcommand{\eqq}{:=}
\newcommand{\eps}{\epsilon}
\newcommand{\veps}{\varepsilon}
\newcommand{\rank}{\text{rank}}
\newcommand{\dims}{\text{dim}}
\newcommand{\diam}{\text{diam}}
\definecolor{orange}{RGB}{160,80,30}
\newcommand{\brown}[1]{{\color{orange} #1}}                    % brown text
%\newcommand{\red}[1]{{\color{red} #1}}
\DeclareMathOperator{\ext}{ext}

\newcommand{\pf}{\noindent {\bf Proof.}\hspace*{0.3cm}}

\newcommand{\epf}{\hfill \rule[-2pt]{6pt}{6pt}}% \\ \vspace*{0.0cm}}

\DeclareMathOperator{\diag}{diag}

%------------------------------------------------------------------------------------
%------------------------------------------------------------------------------------
\newcommand{\Iscr}{\mathcal{I}}
\newcommand{\Fscr}{\mathcal{F}}
\newcommand{\Jscr}{\mathcal{J}}
\newcommand{\Cscr}{\mathcal{C}}
\newcommand{\Rscr}{\mathcal{R}}
\newcommand{\Sscr}{\mathcal{S}}
\newcommand{\Dscr}{\mathcal{D}}
\newcommand{\Pscr}{\mathcal{P}}
\newcommand{\Yscr}{\mathcal{Y}}
\newcommand{\Hbm}{\bm{H}}
\newcommand{\hbm}{\bm{h}}
\newcommand{\dbm}{\bm{d}}
\newcommand{\pbm}{\bm{p}}
\newcommand{\Fbm}{\bm{F}}
\newcommand{\Lbm}{\bm{L}}
\newcommand{\fbm}{\bm{f}}
\newcommand{\cbm}{\bm{c}}
\newcommand{\wbm}{\bm{w}}
\newcommand{\rbm}{\bm{\rho}}
\newcommand{\lbm}{\bm{\lambda}}
\newcommand{\gbm}{\bm{\gamma}}
\newcommand{\abm}{\bm{\alpha}}
\newcommand{\ftcost}{C_{FT}}
\newcommand{\disc}[1]{e^{-r #1}}
\newcommand{\union}[1]{\bigcup_{#1}}
\newcommand{\task}{i} % index for a task
\newcommand{\Stask}{\Theta} % set of tasks
\newcommand{\Admiset}{\mathcal{A}(\Stask)} % set of admissible taks
\newcommand{\profit}{\pi} %profit function
\newcommand{\profitrate}{z}
\newcommand{\fttime}{\tau} % generic fine-tuning time
\newcommand{\State}{\mathcal{S}}

\newcommand{\GBC}{GBC}
\newcommand{\GBCzs}{$\mbox{GBC}_{zs}$}
\newcommand{\GBClc}{$\mbox{GBC}_{lc}$}
%-------------------------------------------

\definecolor{pink}{RGB}{255, 105, 180} % puedes ajustar este color
\definecolor{peach}{RGB}{246, 161, 127} % puedes ajustar este color
% puedes ajustar este color

\newcommand{\natalia}[1]{\textcolor{pink}{#1}}
\newcommand{\denis}[1]{\textcolor{peach}{#1}}
\newcommand{\francisco}[1]{\textcolor{magenta}{#1}}

%------------------------------------------------------------------------------------
\begin{document}

\RUNAUTHOR{Natalia Trigo}


\RUNTITLE{summer paper}

\TITLE{Summer paper II title}


\ARTICLEAUTHORS{

%\AUTHOR{Denis Saur\'e}
%\AFF{Department of Industrial Engineering, University of Chile, Santiago, Chile, \EMAIL{dsaure@dii.uchile.cl}}
\AUTHOR{Natalia Trigo}
\AFF{Anderson UCLA, \EMAIL{natalia.trigo.phd@anderson.ucla.edu}}
}


\ABSTRACT{}

%\KEYWORDS{Sequential $<$ Decision analysis, Combinatorics $<$ Mathematics, Tactics/strategy $<$ Military}


\maketitle

\section{Model Formulation}

Consider an infinite-horizon discrete-time setting, indexed by $t \in \mathbb{N}^+ = \{1, 2, \dots\}$. The environment consists of a single user and two sellers, denoted by $A$ and $B$. 

At each period $t$, the user selects a seller, denoted by the action $y_t \in \{A, B\}$. The user's realized utility in period $t$ is a binary random variable $u_t \in \{0, 1\}$. 

If the user selects Seller B ($y_t = B$), they are offered a single, transparent product. The realized utility is drawn from a Bernoulli distribution with a strictly known parameter $p_0 \in (0, 1)$:
\begin{equation}
    \mathbb{P}(u_t = 1 \mid y_t = B) = p_0.
\end{equation}

If the user selects Seller A ($y_t = A$), Seller A dynamically assigns a product from a finite set $\mathcal{X} = \{1, 2\}$. Let $x_t \in \mathcal{X}$ denote Seller A's choice. The assigned product determines both the probability of a successful outcome for the user and the cost incurred by Seller A. Specifically, product $i \in \mathcal{X}$ yields $u_t = 1$ with probability $p_i \in (0, 1)$ and imposes a cost $c_i \in \mathbb{R}_+$ on Seller A. 
\begin{equation}
    \mathbb{P}(u_t = 1 \mid y_t = A, x_t = i) = p_i.
\end{equation}
Without loss of generality, we assume $c_1 < c_2$. Furthermore, we assume that whenever Seller A is chosen ($y_t = A$), they receive a fixed revenue $R > 0$. The immediate profit for Seller A in period $t$ is thus:
\begin{equation}
    \pi_t^A = 
    \begin{cases} 
      R - c_{x_t} & \text{if } y_t = A, \\
      0 & \text{if } y_t = B.
    \end{cases}
\end{equation}

\section{Information Structure and Learning Dynamics}

The model's core tension stems from an information asymmetry between the user and Seller A. 

\begin{assumption}[User's Misspecified Belief]
The user observes the realized utility $u_t$ but cannot observe the specific product $x_t \in \mathcal{X}$ assigned by Seller A. Consequently, the user operates under the misspecified belief that Seller A offers a single product characterized by an unknown success probability $\theta \in [0, 1]$.
\end{assumption}

The user learns $\theta$ through Bayesian updating. Let the user's prior belief over $\theta$ at $t=1$ be uniformly distributed, given by $\theta \sim \text{Beta}(1, 1)$. Let $S_t$ and $F_t$ denote the cumulative number of successes and failures, respectively, generated by Seller A up to the beginning of period $t$:
\begin{align}
    S_t &= \sum_{\tau=1}^{t-1} \Ind{y_\tau = A, u_\tau = 1}, \\
    F_t &= \sum_{\tau=1}^{t-1} \Ind{y_\tau = A, u_\tau = 0},
\end{align}
The user's posterior belief at period $t$ follows a Beta distribution:
\begin{equation}
    \theta \mid S_t, F_t \sim \text{Beta}(1 + S_t, 1 + F_t).
\end{equation}

The user follows a Thompson Sampling heuristic for their decision policy. At the beginning of period $t$, the user samples an index $\tilde{\theta}_t \sim \text{Beta}(1 + S_t, 1 + F_t)$. Because the expected utility of Seller B ($p_0$) is perfectly known, the user selects Seller A if and only if the sampled index exceeds $p_0$:
\begin{equation}
    y_t = 
    \begin{cases} 
      A & \text{if } \tilde{\theta}_t \geq p_0, \\ 
      B & \text{if } \tilde{\theta}_t < p_0.
    \end{cases}
\end{equation}

\section{Seller A's Optimization Problem}

Seller A observes the sequence of the user's choices and the realized utilities. Therefore, Seller A perfectly tracks $S_t$ and $F_t$, thereby possessing complete information regarding the user's internal belief state.

Let the state of the system from Seller A's perspective be defined by the user's sufficient statistics, $\bm{s}_t = (S_t, F_t) \in \mathbb{N}_0 \times \mathbb{N}_0$. Given state $\bm{s}_t$, the probability that the user chooses Seller A in period $t$, denoted $\rho(S_t, F_t)$, is the complementary cumulative distribution function (CCDF) of the Beta distribution evaluated at $p_0$:
\begin{equation}
    \rho(S_t, F_t) \triangleq \mathbb{P}(\tilde{\theta}_t \geq p_0 \mid S_t, F_t) = \int_{p_0}^{1} \frac{\theta^{S_t} (1-\theta)^{F_t}}{B(1+S_t, 1+F_t)} \, d\theta,
\end{equation}
where $B(\cdot, \cdot)$ is the Beta function.

Seller A faces a Markov Decision Process (MDP) where they must balance immediate cost minimization against the strategic manipulation of the user's belief state to ensure future engagement. Let $\gamma \in (0, 1)$ denote Seller A's discount factor. 

\begin{problem}[Seller A's MDP]
Seller A seeks a stationary policy $\pi_A : \mathbb{N}_0 \times \mathbb{N}_0 \to \mathcal{X}$ that maximizes the expected discounted cumulative profit:
\begin{equation}
    \max_{\pi_A} \mathbb{E}_{\pi_A} \left[ \sum_{t=1}^{\infty} \gamma^{t-1} \Ind{y_t = A} (R - c_{x_t}) \right].
\end{equation}
\end{problem}

Let $V(S, F)$ denote the optimal value function for Seller A at state $(S, F)$. The state transitions occur only when the user selects Seller A. If $y_t = B$, the state remains $(S, F)$ and no profit is accrued, though time advances. The Bellman optimality equation is therefore given by:
\begin{equation}
\begin{aligned}
    V(S, F) =& \big[1 - \rho(S, F)\big] \gamma V(S, F) \\
    &+ \rho(S, F) \max_{x \in \{1, 2\}} \Big\{ (R - c_x) + \gamma \big[ p_x V(S+1, F) + (1 - p_x) V(S, F+1) \big] \Big\}.
\end{aligned}
\end{equation}

By algebraic manipulation, we can isolate $V(S, F)$ to yield a computationally tractable form that resolves the self-referential term arising from periods where Seller A is not chosen:
\begin{equation}
    V(S, F) = \frac{\rho(S, F)}{1 - \gamma + \gamma \rho(S, F)} \max_{x \in \{1, 2\}} \Big\{ (R - c_x) + \gamma \big[ p_x V(S+1, F) + (1 - p_x) V(S, F+1) \big] \Big\}.
\end{equation}

%\end{document}
% \section{Model}

% We consider a dynamic setting with competing LLM providers. The key feature of the model is that the provider chooses which product to use, but the user does not observe this product choice. Instead, the user only observes the realized utility from interacting with the chosen provider. Therefore, the user learns about each provider's reputation, rather than directly learning which product was used.

% There are two providers, indexed by $j \in {A,B}$. Each provider has access to two products, indexed by $k \in {1,2}$. Product $1$ is cheaper for the provider, while product $2$ is more expensive. Let $c_k$ denote the cost of product $k$, with $c_1 < c_2$

% If provider $j$ uses product $k$, the expected utility delivered to the user is $\theta_{j,k}$. We assume that product $2$ generates higher expected utility than product $1$, that is $\theta_{j,2} > \theta_{j,1}$. Thus, product $2$ is better for the user but more costly for the provider.

% \subsection{Hidden Product Choice}

% In each period $t=0,1,2,\ldots$, each provider privately chooses which product to use. Let $a_{j,t} \in \{1,2\}$ denote the product chosen by provider $j$ in period $t$. The product choice $a_{j,t}$ is not observed by the user. The user only observes the realized utility after choosing a provider.

% If the user chooses provider $j$ in period $t$, the realized utility is
% \begin{equation}
% y_t = \theta_{j,a_{j,t}} + \varepsilon_t,
% \end{equation}
% where $\varepsilon_t \sim N(0,\sigma^2)$.

% \subsection{User Beliefs}

% Because the product choice is hidden, the user does not learn separately about $\theta_{j,1}$ and $\theta_{j,2}$. Instead, the user learns about the average utility delivered by provider $j$. Let $\mu_j$ denote the user's perceived average utility from provider $j$.

% The user begins with a prior belief
% \begin{equation}
% \mu_j \sim N(m_{j,0},\tau_{j,0}^2).
% \end{equation}

% At time $t$, the user's belief about provider $j$ is summarized by the posterior mean $m_{j,t}$ and posterior variance $\tau_{j,t}^2$. The belief state is

% \begin{equation}
% b_t =
% ({m_{A,t},\tau_{A,t}^2,m_{B,t},\tau_{B,t}^2})
% \end{equation}

% The posterior mean $m_{j,t}$ represents the user's current estimate of the expected utility delivered by provider $j$. The posterior variance $\tau_{j,t}^2$ represents the user's uncertainty about provider $j$.

% When the user chooses provider $j$ and observes $y_t$, the user updates the belief about provider $j$ while for the other provider remains unchanged.

% \subsection{Timing}

% The timing in each period is as follows.

% \begin{enumerate}
% \item Each provider $j$ privately chooses a product $a_{j,t} \in {1,2}$.
% \item The user observes the belief state $b_t$, but does not observe the product choices.
% \item The user chooses one provider $j_t \in {A,B}$.
% \item The chosen provider delivers noisy utility $y_t$.
% \item The user observes $y_t$, but not the product used.
% \item The user updates the belief about the chosen provider.
% \end{enumerate}

% This timing captures a hidden-action environment. The provider's product choice affects the user's realized utility, but the user cannot directly observe the action that generated that utility.

% \subsection{Myopic User}

% We first consider a myopic user. A myopic user chooses the provider with the highest current posterior mean and does not account for the value of learning.

% Given belief state $b_t$, the myopic user chooses
% \begin{equation}
% j_t^M
% \in
% \arg\max_{j \in {A,B}}
% m_{j,t}.
% \end{equation}

% Thus, the myopic user chooses the provider that currently appears to deliver the highest average utility. Since the user cannot observe product choice, the user does not choose based on whether the provider is using product $1$ or product $2$. The user only chooses based on provider reputation.

% \subsection{Forward-Looking User}

% We also consider a forward-looking user. A forward-looking user recognizes that choosing a provider today produces information that may improve future choices. Let $V(b_t)$ denote the user's continuation value when the belief state is $b_t$. The user's dynamic problem is

% \begin{equation}
% V(b_t)=
% \max_{j \in {A,B}}
% \mathbb{E}
% \left[
% y_t
% +
% \beta V(b_{t+1})
% \mid b_t,j
% \right],
% \end{equation}

% where $\beta \in (0,1)$ is the user's discount factor. The expectation is taken over the realized utility $y_t$ and the resulting posterior belief $b_{t+1}$.

% The forward-looking user may choose a provider with a lower posterior mean if that provider has sufficiently high uncertainty. Therefore, the forward-looking user trades off current expected utility and the value of information.

% \subsection{Provider Payoff}

% If provider $j$ is chosen in period $t$, it receives revenue $R$ and pays the cost of the product it privately selected. Let $x_{j,t}$ be an indicator for whether provider $j$ is chosen. The one-period profit of provider $j$ is
% \begin{equation}
% \pi_{j,t}
% =x_{j,t}
% \left(R-c_{a_{j,t}}\right).
% \end{equation}

% If provider $j$ is not chosen, then $x_{j,t}=0$ and the provider earns zero profit in that period.

% \subsection{Provider Problem}

% Provider $j$ chooses which product to use in order to maximize expected discounted profit. Let $W_j(b_t)$ denote provider $j$'s continuation value when the user's belief state is $b_t$. Then provider $j$ solves
% \begin{equation}
% W_j(b_t)
% =
% x_j(b_t)
% \max_{a_{j,t}\in\{1,2\}}
% \mathbb{E}
% \left[
% R-c_{a_{j,t}}+\delta W_j(b_{t+1})
% \mid b_t, x_{j,t}=1, a_{j,t}
% \right]
% +
% \bigl(1-x_j(b_t)\bigr)\delta W_j(b_t).
% \end{equation}

% where $\delta \in (0,1)$ is the provider's discount factor.

% The provider's action $a_{j,t}$ affects current profit through the cost $c_{a_{j,t}}$. It also affects future profit through the user's belief update. If the provider chooses product $2$, it pays a higher cost today, but it generates a higher expected utility signal. This can improve the user's posterior belief and increase the probability that the provider is chosen in the future. If the provider chooses product $1$, it saves cost today, but it generates a lower expected utility signal and may damage future reputation.

% Therefore, the provider faces a reputation trade-off. Product $2$ is an investment in reputation, while product $1$ is a cost-saving action that may reduce future demand.

% \subsection{Provider problem when he doesn't know the belief state}

% We now consider a setting in which providers do not observe the user's belief state. In particular, provider $j$ does not observe $m_{A,t}$, $\tau^2_{A,t}$, $m_{B,t}$, or $\tau^2_{B,t}$. Instead, provider $j$ only observes its own demand history.

% Let
% \begin{equation}
% x_{j,t}
% =
% \begin{cases}
% 1 & \text{if the user chooses provider } j \text{ in period } t,\\
% 0 & \text{otherwise.}
% \end{cases}
% \end{equation}

% The information available to provider $j$ at the beginning of period $t$ is
% \begin{equation}
% h^j_t = (x_{j,0},x_{j,1},\dots,x_{j,t-1}).
% \end{equation}


% Therefore, provider $j$ cannot condition its product choice directly on the user's posterior beliefs. Instead, it must infer its reputation from the user's past choices. A natural summary of this history is provider $j$'s recent market share,

% \begin{equation}
%     \frac{1}{H}\sum_{\ell=t-H}^{t-1} x_{j,\ell}
% \end{equation}

% where $H$ is the length of the provider's memory window. The statistic $r_{j,t}$ captures how frequently the user has recently selected provider $j$.

% In this regime, provider $j$ chooses a product according to a policy

% \begin{equation}
%     a_{j,t} = \pi_j(r_{j,t}),
% \end{equation}


% rather than a policy that depends on the full belief state. For example, provider $j$ may use the high-quality product when recent demand is low or intermediate, and use the low-cost product when recent demand is high. One simple policy is

% \begin{equation}
%     \begin{cases}
% 2 & \text{if } r_{j,t}<\bar r,\\
% 1 & \text{if } r_{j,t}\geq \bar r,
% \end{cases}
% \end{equation}

% where $\bar r$ is a demand threshold.

% The economic interpretation is that recent user choices serve as a noisy signal of reputation. If provider $j$ is chosen frequently, it infers that the user likely has a favorable belief about it, so it may save costs by using product 1. If provider $j$ is chosen infrequently, it infers that its reputation may be weak, so it may invest in product 2 in order to generate better signals when it is chosen.

% More generally, the provider's problem can be written as

% \begin{equation}
%     \max_{a_{j,t}\in{1,2}}
% \mathbb{E}
% \left[
% x_{j,t}(R-c_{a_{j,t}})
% +
% \delta W_j(h^j_{t+1})
% \mid h^j_t,a_{j,t}
% \right].
% \end{equation}

% The expectation is taken over the user's choice, the realized utility signal, and the evolution of future choice histories. The key difference from the full-information case is that the provider does not know the belief state that drives the user's choice. It only observes the realized sequence of choices.

% This information structure creates an additional role for user exploration. If the user is myopic, a provider that falls behind may stop being chosen, which limits its opportunity to repair its reputation. If the user uses Thompson Sampling, the losing provider may still be chosen occasionally because of exploration. These exploratory choices give the provider opportunities to use the high-quality product and improve the user's future beliefs. Thus, Thompson Sampling affects provider incentives not only by changing user learning, but also by giving providers more demand feedback.

\end{document}